\chapter{Hemorrhoids}

\section*{Definition}
Hemorrhoids are symptomatic enlargement and distal displacement of the normal anal cushions (vascular tissue) located above (internal) or below (external) the dentate line. They are graded I--IV based on degree of prolapse: I (no prolapse), II (prolapse with defecation, spontaneous reduction), III (requires manual reduction), IV (irreducible).

\section*{Pathophysiology}
Weakening of the supporting connective tissue within the anal cushions allows the hemorrhoidal vascular plexus to prolapse distally. Straining during defecation, increased intra-abdominal pressure, and venous congestion cause engorgement, bleeding, and thrombosis of the hemorrhoidal vessels.

\section*{Causes}
\begin{itemize}
  \item \textbf{Constipation and straining:} Chronic constipation with prolonged Valsalva (most common cause)
  \item \textbf{Pregnancy:} Increased intra-abdominal pressure, hormonal changes, and venous congestion
  \item \textbf{Dietary:} Low-fiber diet, inadequate fluid intake, excessive caffeine/alcohol
  \item \textbf{Lifestyle:} Prolonged sitting, obesity, sedentary lifestyle, occupations with heavy lifting
  \item \textbf{Other:} Portal hypertension (cirrhosis), pelvic tumors, anal intercourse, family predisposition
\end{itemize}

\section*{When to See a Doctor}
See your doctor if:
\begin{itemize}
  \item Painful external swelling with acute onset (suggests thrombosed external hemorrhoid)
  \item Bright red blood in stool or on toilet paper (must exclude other causes of hematochezia)
  \item Prolapsing tissue requiring manual reduction or irreducible
  \item Persistent symptoms despite lifestyle and dietary modifications
  \item Change in bowel habits, weight loss, or abdominal pain suggesting colorectal pathology
\end{itemize}

\section*{Related Symptoms}
\begin{itemize}
  \item Painless rectal bleeding, anal pain, pruritus ani, feeling of incomplete evacuation, prolapsing tissue, perianal swelling
\end{itemize}
\textbf{Important:} Internal hemorrhoids above the dentate line are insensate and typically present with painless bleeding, while external hemorrhoids below the dentate line are innervated and become painful with thrombosis.

\section*{Self-Care / Home Management}
\begin{itemize}
  \item Increase dietary fiber to 25--35 g/day (psyllium husk, methylcellulose) and ensure adequate water intake (8--10 cups/day)
  \item Sitz baths (warm water) for 10--15 minutes after each bowel movement and 2--3 times daily
  \item Topical hydrocortisone 1\% cream or witch hazel pads (Tucks) for pruritus and inflammation (limit to 7 days)
  \item Stool softeners (docusate) to prevent straining; avoid prolonged sitting on the toilet
  \item Cold packs for acute external thrombosed hemorrhoids to reduce edema and pain
\end{itemize}

\section*{How Your Doctor Will Evaluate You}
\begin{itemize}
  \item Visual inspection for external hemorrhoids, skin tags, prolapse
  \item Digital rectal exam to assess internal hemorrhoids and exclude other mass lesions
  \item Anoscopy: gold standard for visualizing internal hemorrhoids (position at 3, 7, and 11 o'clock typically)
  \item Colonoscopy if alarm features (age over 50, weight loss, anemia, family history of colon cancer, change in caliber) are present
  \item Reducibility assessment and grading (I--IV)
\end{itemize}

\section*{Treatment Options}
\begin{itemize}
  \item Grade I--II: Dietary modification, increased fiber, warm water irrigation (toilet seat bidet), topical therapy
  \item Grade II--III: Office-based procedures including rubber band ligation, sclerotherapy, or infrared coagulation
  \item Thrombosed external hemorrhoid: Excision under local anesthesia within 72 hours for faster relief vs. conservative management
  \item Grade III--IV: Hemorrhoidectomy (excisional) or stapled hemorrhoidopexy for symptomatic prolapse
  \item Doppler-guided hemorrhoidal artery ligation (DG-HAL) as alternative surgical option with less post-operative pain
\end{itemize}

\clearpage
